\subsection{完全割理想的刻画}

设 A 为环，J 为 A 的理想。分次 A-模 \(\bigoplus_{n\in\mathbf{N}}J^{n}\) 具有一个自然的分次 A-代数结构，它由环 A 中的乘法导出；于是 J 到 \(J^{1}\) 的恒同映射扩张为分次 A-代数的一个满射同态，称为典范同态

\[
\alpha_ {J}: \mathsf {S} _ {A} (J) \longrightarrow \bigoplus_ {n \in N} J ^ {n}.
\]

通过向环 A/J 作纯量扩张，由 \(\alpha_{\mathrm{J}}\) 导出一个分次 A/J-代数的满射同态，也称为典范同态

\[
\beta_ {J}: \mathsf {S} _ {\mathrm{A} / J} (\mathrm{J} / \mathrm{J} ^ {2}) \longrightarrow \operatorname{gr} _ {J} (\mathrm{A}),
\]

\(\operatorname{avec} \mathrm{gr}_{\mathrm{J}}(\mathrm{A}) = \bigoplus_{n \in \mathbf{N}} \mathrm{J}^n / \mathrm{J}^{n+1}\) 。

定理 1.--- 设 A 为 Noether 环，J 为 A 的理想。下列条件等价：

\begin{enumerate}
\def\labelenumi{(\roman{enumi})}
\item
  理想 J 是完全割的；
\item
  理想 J 在每个极大理想 \(\mathfrak{m} \in V(J)\) 处都完全割；
\item
  A/J-模 J/J² 是投射的，且典范同态 \(\alpha_{J}: S_{\Lambda}(J) \longrightarrow \bigoplus_{n \in \mathbb{N}} J^{n}\) 是双射；
\item
  \(\Lambda/\mathrm{J}\) -模 \(\mathrm{J}/\mathrm{J}^2\) 是投射的，且典范同态 \(\beta_{\mathrm{J}}:\mathbf{S}_{\mathrm{A}/\mathrm{J}}(\mathrm{J}/\mathrm{J}^2)\longrightarrow\mathrm{gr}_{\mathrm{J}}(\mathrm{A})\) 是双射。
\end{enumerate}

(i)⇒(ii)：这是显然的。

(ii)⇒(iii)：设条件 (ii) 满足。只需证明：对 A 的每个极大理想 m，\(A_{m}/J_{m}\) -模 \(J_{m}/J_{m}^{2}\) 是自由的，且同态 \(\alpha_{J_{m}}: S_{A_{m}}(J_{m}) \longrightarrow \bigoplus_{n} J_{m}^{n}\) 是双射 (II, § 5, 第 2 节, 定理 1 及 § 3, 第 3 节, 定理 1)。但当 m 不属于 V(J) 时这些断言是显然的，因为此时有 \(J_{m} = A_{m}\)；而当 m 属于 V(J) 时，它们由 A, X, 第 160 页的定理 1 及第 161 页的注得出。

\begin{enumerate}
\def\labelenumi{(\roman{enumi})}
\setcounter{enumi}{2}
\tightlist
\item
  \(\Rightarrow\) (iv)：这是显然的。
\end{enumerate}

\((\mathrm{iv}) \Rightarrow (\mathrm{i})\)：设条件 (iv) 满足；设 \(\mathfrak{p}\) 为 A 的一个包含 J 的素理想。则 \(A_{\mathfrak{p}} / J_{\mathfrak{p}}\) -模 \(J_{\mathfrak{p}} / J_{\mathfrak{p}}^2\) 是自由的。设 \(\mathbf{x} = (x_1, \ldots, x_r)\) 为 \(J_{\mathfrak{p}}\) 的元素的一个序列，它提升 \(J_{\mathfrak{p}} / J_{\mathfrak{p}}^2\) 的一组基。序列 \(\mathbf{x}\) 生成 \(J_{\mathfrak{p}}\)（Nakayama 引理），且按构造满足命题 1 的条件 (iv)。于是理想 \(J_{\mathfrak{p}}\)（\(A_{\mathfrak{p}}\) 中的）是完全割的，即 J 在 \(\mathfrak{p}\) 处完全割。

注 1. 设理想 J 是完全割的；设 \((x_1, \ldots, x_r)\) 为 J 的元素的一个序列，使得对每个极大理想 \(\mathfrak{m} \in V(J)\)，\(x_i\) 在 J/mJ 中的典范像构成这个 A/m-向量空间的一组基。则 A/J-模 \(J/J^2\) 是自由的，且在 \(J/J^2\) 中 \(x_i\) 的典范像构成它的一组基：事实上只需验证：\(x_i\) 的像构成 \(A_{\mathfrak{m}}/J_{\mathfrak{m}}\) -模 \(J_{\mathfrak{m}}/J_{\mathfrak{m}}^2\) 的一组基（对每个 \(\mathfrak{m} \in V(J)\)） (II, § 3, 第 3 节, 定理 1)，而这由同上引文第 2 节的命题 5 与推论 2 得出，因为 \(A_{\mathfrak{m}}/J_{\mathfrak{m}}\) -模 \(J_{\mathfrak{m}}/J_{\mathfrak{m}}^2\) 是投射的 (定理 1)。

推论 1.--- 设 \(\widehat{\mathrm{A}}\) 为 A 关于 J-进拓扑的分离完备化，\(\widehat{\mathrm{J}} = \mathrm{J}\widehat{\mathrm{A}}\) 为 J 的分离完备化。为使理想 \(\widehat{\mathrm{J}}\)（\(\widehat{\mathrm{A}}\) 中的）完全割，必须且只需 A 的理想 J 完全割。

事实上，典范映射 \(\mathrm{gr}_{\mathrm{J}}(\mathsf{A}) \to \mathrm{gr}_{\widehat{\mathrm{J}}}(\widehat{\mathsf{A}})\) 是分次环的同构；因此只需应用判别法 (iv)。

更一般地：

推论 2.--- 设 ρ : A → B 为一个 Noether 环的同态，它使 B 成为平坦 A-模，并设 J 为 A 的理想。

\begin{enumerate}
\def\labelenumi{\alph{enumi})}
\item
  若 J 完全割，则 B 的理想 JB 完全割。
\item
  设 B 的理想 JB 完全割，且 \(\mathfrak{m} \in \mathrm{V}(\mathrm{J})\) 的每个极大理想都是 B 的某个极大理想的原像。则理想 J 完全割。例如当 B 是忠实平坦 A-模时就是这种情形。
\end{enumerate}

首先注意，由于 B 是平坦 A-模，\(J^n \otimes_A B\) 等同于 \(J^n B\)，且 \((J^n / J^{n+1}) \otimes_{A/J} (B/JB)\) 等同于 \(J^n B / J^{n+1} B\)（对每个整数 \(n \geqslant 0\)）。于是断言 a) 由判别法 (iii) 得出。在 b) 的假设下，A/J-模 B/JB 是忠实平坦的 (I, § 2, 第 7 节, 命题 8 的推论 2 及 § 3, 第 5 节, 命题 9)，且考虑到 I, § 3, 第 1 节, 命题 2 及第 6 节, 命题 12，J 依判别法 (iv) 完全割。最后的断言由 I, § 3, 第 5 节的命题 8 得出。

推论 3.-- 设 A 为 Noether 环，J 为 A 的完全割理想。若 A 是 Macaulay 环（相应地 Gorenstein 环），则 A/J 亦然。

事实上，设 m 为 A 的一个包含 J 的极大理想。理想 \(J_{m}\)（\(\Lambda_{m}\) 中的）由一个 \(A_{m}\) -正则序列生成；于是 \((A/J)_{m}\) 是 Macaulay 环（相应地 Gorenstein 环），依据是 § 2, 第 5 节的例 6（相应地 § 3, 第 9 节的例 2）。

注 2. — 称 Noether 环 A 为完全交环，如果对 A 的每个极大理想 m，完备局部环 \(\widehat{A_{m}}\) 同构于一个正则完备 Noether 局部环对一个完全交理想的商。由上面的推论 3 及 § 3, 第 8 节的命题 12 的推论 2 得出，这样的环是 Gorenstein 环。

